% 2_assumptions.tex — 模型假设

为了建立合理且可求解的数学模型，本文做出以下基本假设：

\begin{enumerate}[label=\textbf{假设\arabic*：}, leftmargin=*]
    \item \textbf{导弹飞行假设}：来袭导弹以恒定速度300 m/s沿直线飞行，方向始终指向假目标中心，不考虑导弹的机动变轨或末端制导修正。
    \textit{合理性}：空地导弹在末制导段前通常保持稳定的飞行弹道；在烟幕干扰的作战窗口内，导弹尚未进入末端机动阶段。

    \item \textbf{无人机飞行假设}：无人机在接到任务后可瞬时调整飞行方向，之后以恒定速度、恒定航向沿直线等高度飞行，不考虑转弯半径和加减速过程。
    \textit{合理性}：无人机的任务调整时间远小于整体作战时间尺度；现代无人机具备快速航向调整能力；假设简化了轨迹优化的复杂度。

    \item \textbf{烟幕云团假设}：烟幕干扰弹起爆后瞬时形成严格的球形烟幕云团，云团内烟幕浓度均匀，对视线遮蔽效果为二值化（完全遮挡或完全不遮挡），云团以恒定速度3 m/s垂直下沉。
    \textit{合理性}：试验数据表明10 m半径范围内烟幕浓度足以提供有效遮蔽，球状模型是烟幕扩散的合理一阶近似\cite{zhang2020}。

    \item \textbf{自由落体假设}：烟幕干扰弹脱离无人机后仅受重力作用，忽略空气阻力对水平运动的影响。
    \textit{合理性}：烟幕干扰弹在短时（$<$20 s）自由落体阶段，空气阻力对水平速度的影响远小于重力对垂直位置的影响；且干扰弹质量较大，风阻效应有限。

    \item \textbf{空域安全假设}：所有烟幕干扰弹的起爆高度不低于地面，且烟幕云团在下沉过程中不会因接触地面而提前消散。
    \textit{合理性}：实际作战中干扰弹通常在数百米以上高度起爆，烟幕云团在20 s有效期内一般不会触地。

    \item \textbf{通信与控制假设}：控制中心与无人机之间的通信延迟可忽略不计，无人机可精确按照指定时刻和位置投放烟幕干扰弹，起爆时刻可通过时间引信精确控制。
    \textit{合理性}：现代军用数据链的通信延迟在毫秒级，相比秒级的作战时间窗口可忽略。

    \item \textbf{真假目标假设}：来袭导弹始终以假目标为瞄准点，在被烟幕遮蔽后不会重新搜索或切换目标。真目标与假目标的空间位置关系固定不变。
    \textit{合理性}：这是烟幕干扰作战的基本前提；导弹在失去目标锁定后重新捕获的概率较低，且已超出本题研究范围。
\end{enumerate}
